\section{Estimación de costes}\label{cap:estimacionCostes}

A continuación se presenta el presupuesto estimado del proyecto, calculado como si este fuese desarrollado por un equipo profesional completo, siguiendo los perfiles habituales en el desarrollo de una aplicación móvil de estas características. Se trata de una estimación teórica con fines de planificación y no de los costes reales asumidos por el autor, quien ha desarrollado el TFG de forma individual.

\subsection{Costes de personal}

Un proyecto de estas características, correspondiente al desarrollo completo de una aplicación móvil sobre la Semana Santa, requeriría en un entorno profesional un equipo multidisciplinar. Se han considerado los siguientes seis perfiles, entre los que se reparten las 300 horas estimadas de duración del proyecto:

\begin{itemize}
    \item \textbf{Jefe de Proyecto}: coordinación general, planificación y seguimiento.
    \item \textbf{Analista Funcional}: recogida y especificación de requisitos, análisis del dominio.
    \item \textbf{Diseñador UX/UI}: diseño de la experiencia e interfaz de usuario.
    \item \textbf{Desarrollador Backend}: implementación del servidor, base de datos y API.
    \item \textbf{Desarrollador Frontend/Mobile}: implementación de la aplicación Android.
    \item \textbf{Tester/QA}: pruebas funcionales, de regresión y control de calidad.
\end{itemize}

Para cada perfil se ha calculado el coste por hora a partir del sueldo bruto anual medio en España, incrementado en un 30\% en concepto de cotizaciones a la Seguridad Social a cargo de la empresa, y dividido entre 1.800 horas anuales de jornada completa. La Tabla~\ref{tab:horas_personal} recoge el reparto de horas y el coste asociado a cada perfil.

\begin{table}[ht!]
\centering
\begin{tabular}{|l|r|r|r|}
\hline
\textbf{Perfil} & \textbf{Horas} & \textbf{Coste/hora} & \textbf{Coste total} \\
\hline
Jefe de Proyecto \cite{salarioJefeProyecto}          & 45  & 33€ & 1.485,00€ \\
\hline
Analista Funcional \cite{salarioAnalista}            & 45  & 27€ & 1.215,00€ \\
\hline
Diseñador UX/UI \cite{salarioUXUI}                   & 30  & 27€ &   810,00€ \\
\hline
Desarrollador Backend \cite{salarioBackend}          & 75  & 26€ & 1.950,00€ \\
\hline
Desarrollador Frontend/Mobile \cite{salarios}        & 75  & 19€ & 1.425,00€ \\
\hline
Tester/QA \cite{salarioTester}                       & 30  & 17€ &   510,00€ \\
\hline
\textbf{Total} & \textbf{300} & & \textbf{7.395,00€} \\
\hline
\end{tabular}
\caption{Reparto de horas y coste de personal por perfil}
\label{tab:horas_personal}
\end{table}

\subsection{Costes de equipo y amortización}

El equipo del proyecto necesitaría un ordenador portátil por persona, además de un dispositivo Android dedicado a las pruebas. Se considera una duración total del proyecto de 5 meses.

\begin{itemize}
    \item \textbf{Ordenadores portátiles}: 6 unidades valoradas en 1.000€ cada una, con una vida útil estimada de 5 años (60 meses). La amortización mensual por unidad es de 16,66€, lo que supone 83,30€ por portátil durante los 5 meses del proyecto, y 500,00€ en total para los 6 puestos de trabajo.
    \item \textbf{Dispositivo Android de pruebas}: 1 unidad valorada en 300€, con una vida útil estimada de 3 años (36 meses). La amortización mensual es de 8,33€, lo que supone 41,65€ para los 5 meses del proyecto.
\end{itemize}

Ahora bien, estos importes corresponden a la ocupación de los equipos durante los 5 meses completos del proyecto. Dado que el trabajo real asciende a 300 horas, frente a una capacidad disponible de 750 horas en ese periodo (150 horas/mes de jornada completa, según el criterio de 1.800 horas anuales aplicado en el apartado anterior), el resto del tiempo los equipos podrían destinarse a otro uso. Por ello, la amortización se imputa de forma proporcional al uso real del proyecto (300/750 = 40\%).

El coste de amortización así calculado asciende a \textbf{216,66€}, tal y como recoge la Tabla~\ref{tab:amortizacion}.

\begin{table}[ht!]
\centering
\begin{tabular}{|l|r|r|r|r|}
\hline
\textbf{Recurso} & \textbf{Valor} & \textbf{Vida útil} & \textbf{Amort. (5 meses)} & \textbf{Amort. proporcional (300h, 40\%)} \\
\hline
Portátiles (x6) & 6.000,00€ & 5 años & 500,00€ & 200,00€ \\
\hline
Dispositivo Android (x1) & 300,00€ & 3 años & 41,65€ & 16,66€ \\
\hline
\textbf{Total amortización} & & & 541,65€ & \textbf{216,66€} \\
\hline
\end{tabular}
\caption{Amortización de los recursos de equipo, prorrateada al uso real del proyecto}
\label{tab:amortizacion}
\end{table}

\subsection{Costes de licencias de software}

Buena parte de las herramientas empleadas disponen de planes gratuitos o licencias universitarias suficientes para un proyecto de este tamaño. Sin embargo, un equipo profesional de 6 personas necesitaría además herramientas colaborativas de pago para la gestión del proyecto y el diseño. La Tabla~\ref{tab:licencias} recoge el coste y la fuente de cada licencia.

\begin{table}[ht!]
\centering
\begin{tabular}{|l|r|l|}
\hline
\textbf{Herramienta} & \textbf{Coste} & \textbf{Fuente} \\
\hline
Android Studio & 0,00€ & \cite{androidstudio} \\
\hline
Visual Studio Code & 0,00€ & \cite{vscode} \\
\hline
Overleaf (plan gratuito) & 0,00€ & --- \\
\hline
Google Firebase (plan gratuito) & 0,00€ & \cite{firebasefcm} \\
\hline
Google Maps API (crédito gratuito) & 0,00€ & \cite{googlemapsapi} \\
\hline
GitHub (plan gratuito) & 0,00€ & \cite{github} \\
\hline
Cuenta de desarrollador Google Play (pago único) & 23,00€ & \cite{googlePlayFee} \\
\hline
Figma Professional (1 licencia, 5 meses) & 69,00€ & \cite{figmaPricing} \\
\hline
Jira Software Standard (6 usuarios, 5 meses) & 218,40€ & \cite{jiraPricing} \\
\hline
\textbf{Total licencias} & \textbf{310,40€} & \\
\hline
\end{tabular}
\caption{Coste de licencias de software}
\label{tab:licencias}
\end{table}

\subsection{Otros costes indirectos}

Un equipo de 6 personas trabajando de forma presencial o híbrida necesitaría un espacio de oficina. Se ha tomado como referencia el alquiler de un despacho de 22 m\textsuperscript{2} en un espacio de coworking en el centro de Valladolid, con todos los suministros incluidos (calefacción, refrigeración, electricidad, wifi y limpieza), por un importe de 475€ al mes \cite{oficinaValladolid}. Para los 5 meses de duración del proyecto, el coste estructural asciende a 2.375,00€.

Al igual que con la amortización de equipos, este importe corresponde a la disponibilidad del espacio durante los 5 meses completos, mientras que el uso real del proyecto es de 300 horas frente a una capacidad de 750 horas en ese periodo (40\%); el resto del tiempo el espacio podría destinarse a otro uso. Aplicando este mismo criterio de proporcionalidad, el coste estructural imputado al proyecto asciende a \textbf{950,00€} (2.375,00€ × 40\%).

\subsection{Coste total}

La Tabla~\ref{tab:costes_totales} resume el presupuesto completo del proyecto, considerando personal, amortización de equipo, licencias de software y costes estructurales —estos dos últimos prorrateados a las 300 horas reales de uso— para una duración de 300 horas repartidas en 5 meses.

\begin{table}[ht!]
\centering
\begin{tabular}{|l|r|}
\hline
\textbf{Elemento} & \textbf{Coste} \\
\hline
Personal (Tabla~\ref{tab:horas_personal}) & 7.395,00€ \\
\hline
Amortización de equipo, prorrateada (Tabla~\ref{tab:amortizacion}) & 216,66€ \\
\hline
Licencias de software (Tabla~\ref{tab:licencias}) & 310,40€ \\
\hline
Costes estructurales, prorrateados (oficina y suministros) & 950,00€ \\
\hline
\textbf{Total} & \textbf{8.872,06€} \\
\hline
\end{tabular}
\caption{Costes totales del proyecto}
\label{tab:costes_totales}
\end{table}

\subsection{Coste estimado de Firebase en producción}\label{cap:costeFirebaseProduccion}

Los costes anteriores corresponden únicamente al desarrollo del TFG, durante el cual el uso de Firebase se ha mantenido dentro de la capa gratuita (plan \textit{Spark}). Sin embargo, tal y como se señala en la Sección~\ref{cap:puntosDebiles}, el paso a producción con usuarios reales implicaría migrar al plan de pago \textit{Blaze}, facturado según el consumo real. A modo orientativo, y tomando las tarifas oficiales vigentes del plan Blaze para Cloud Firestore \cite{firebasepricing} --0,06 € por cada 100.000 lecturas, 0,18 € por cada 100.000 escrituras y 0,26 € por GB almacenado al mes, descontando la cuota gratuita diaria de 50.000 lecturas y 20.000 escrituras--, se puede estimar el coste de una semana de Semana Santa en una única ciudad piloto:

\begin{itemize}
    \item \textbf{Lecturas}: suponiendo 3.000 usuarios simultáneos consultando el mapa en tiempo real durante una media de 3 horas al día, con actualizaciones cada 30 segundos (RNF-08), se generarían unas 360 lecturas por usuario y día, es decir, aproximadamente 1.080.000 lecturas diarias (7.560.000 en los 7 días de Semana Santa). Descontando la cuota gratuita, el coste asciende a unos \textbf{4,30 €}.
    \item \textbf{Escrituras}: suponiendo 200 cofrades compartiendo su ubicación durante 2 horas al día (240 escrituras por cofrade y día), se generarían unas 48.000 escrituras diarias (336.000 en los 7 días). Descontando la cuota gratuita, el coste asciende a unos \textbf{0,35 €}.
    \item \textbf{Almacenamiento y notificaciones}: el volumen de datos almacenado (perfiles, códigos de acceso, información de cofradías y procesiones) se mantendría muy por debajo de 1~GB, dentro de la cuota gratuita; Cloud Messaging no tiene coste en el plan Blaze \cite{firebasepricing}.
\end{itemize}

Bajo estas hipótesis, el coste de Firebase para una semana de Semana Santa en una ciudad piloto rondaría los \textbf{5 €}, una cifra reducida gracias al bajo coste unitario de las operaciones de Firestore. No obstante, este importe crecería proporcionalmente al escalar la aplicación a más ciudades y usuarios, tal y como se plantea en el trabajo futuro (Sección~\ref{cap:trabajo_futuro}), por lo que esta estimación --y las hipótesis de partida sobre el número de usuarios y ciudades-- debería revisarse y ajustarse junto con el tutor antes de un despliegue real en producción.